\chapter{Pruebas y verificación}\label{cap:pruebas}

Este capítulo describe cómo se comprueba que el sistema hace lo que debe. Presenta la
estrategia de pruebas, los gates de evidencia automáticos, la verificación independiente del
objetivo de cada fase y la revisión de código de todo el repositorio, con sus resultados en
cifras reales tomadas de \texttt{docs/verification/} y de los informes de verificación de las
fases.

\section{Estrategia de pruebas}\label{sec:pru-estrategia}

La estrategia combina cuatro niveles. Las pruebas unitarias y de integración del backend se
ejecutan con \texttt{pytest} contra una instancia real de PostgreSQL 18.6, sin sustituirla
por una base en memoria, porque el sistema depende de restricciones, transacciones y búsqueda
de trigramas que solo se comportan igual sobre el motor real. Las pruebas del frontend usan
\texttt{vitest} para la lógica y los componentes, con una prueba de paridad de claves de
traducción. Las pruebas de extremo a extremo usan Playwright sobre un navegador real, con la
biblioteca \texttt{axe} para las comprobaciones automáticas de accesibilidad frente a las
Pautas de Accesibilidad para el Contenido Web (WCAG) 2.2~\cite{wcag22}. Por último, una
prueba de humo se ejecuta contra la URL pública desplegada y rechaza un despliegue con la
revisión equivocada, tráfico entre orígenes inesperado o peticiones fallidas.

La Tabla~\ref{tab:suites} recoge los resultados de las suites en el momento del cierre de la
Fase 01.1.

\begin{table}[htp]
\centering
\begin{tabular}{ | l | l | l | }
    \hline
    \textbf{Suite} & \textbf{Herramienta} & \textbf{Resultado} \\
    \hline
    Backend & \texttt{pytest} sobre PostgreSQL & 209 pruebas superadas \\
    \hline
    Recomendaciones & \texttt{pytest} & 21 pruebas superadas \\
    \hline
    Frontend, unidad y paridad & \texttt{vitest} & 16 pruebas superadas \\
    \hline
    Construcción del frontend & \texttt{next build} & correcta, sin errores de tipos \\
    \hline
    Extremo a extremo y accesibilidad & Playwright y \texttt{axe} & 41 pruebas superadas; cero incidencias críticas o graves de \texttt{axe} \\
    \hline
\end{tabular}
\caption{Resultados de las suites de pruebas al cierre de la Fase 01.1.}
\label{tab:suites}
\end{table}

Las comprobaciones de \texttt{axe} cubren el catálogo, el detalle, el registro, las
recomendaciones, la página de inicio, el acceso y las fuentes, en escritorio a 1280 por 800
píxeles y en móvil a 375 por 812 píxeles.

\section{Gates de evidencia}\label{sec:pru-gates}

Además de las pruebas, el repositorio tiene gates deterministas que protegen la evidencia. El
gate \texttt{check-secrets.ps1} escanea el árbol de Git, los paquetes de construcción, las
capas de las imágenes y los registros capturados en busca de cadenas con forma de credencial,
y se autoverifica primero contra señuelos sintéticos. El gate \texttt{check-evidence.ps1}
valida el esquema, los tipos, los actores, las rutas, los hashes y la cobertura del registro
de agentes, probando primero con entradas inválidas. El gate \texttt{check-dependencies.ps1}
comprueba que toda dependencia directa está en la lista permitida, fijada a versión exacta,
con su fichero de bloqueo. Los gates \texttt{verify-igdb-*.ps1} comprueban el contenido del
registro de decisión de IGDB, la evidencia del sondeo y la ejecución de importación sobre una
base de datos fresca.

\section{Verificación del objetivo de fase}\label{sec:pru-verificacion-fase}

Cada fase tiene una verificación independiente de su objetivo, realizada por un subagente
distinto del que implementó el trabajo, que comprueba el código contra los criterios de éxito
de la hoja de ruta en lugar de leer los resúmenes de plan. La Fase 1 se verificó con cinco de
cinco criterios de éxito, con dos elementos manuales de accesibilidad aceptados como
limitaciones nombradas: el reflujo al cuatrocientos por ciento de zoom y una pasada con
lector de pantalla. La Fase 01.1 se verificó con cinco de cinco criterios, tras una
reverificación del criterio de cuentas simuladas plurales que al cierre inicial estaba
parcial. La Tabla~\ref{tab:criterios} resume los criterios de la Fase 01.1.

\begin{table}[htp]
\centering
\begin{tabular}{ | c | p{10.5cm} | c | }
    \hline
    \textbf{Criterio} & \textbf{Enunciado} & \textbf{Veredicto} \\
    \hline
    SC1 & Catálogo a escala real con fuente IGDB, con licencia, atribución y procedencia documentadas. & Verificado \\
    \hline
    SC2 & Ordenar y filtrar el catálogo por plataforma, género y otros metadatos. & Verificado \\
    \hline
    SC3 & Varias cuentas precargadas distintas inician y cierran sesión de forma independiente. & Verificado \\
    \hline
    SC4 & La interfaz se lee como un producto de catalogación profesional. & Verificado \\
    \hline
    SC5 & Página de recomendaciones por géneros según la actividad propia, distinta del baseline de popularidad. & Verificado \\
    \hline
\end{tabular}
\caption{Criterios de éxito de la Fase 01.1 y su veredicto.}
\label{tab:criterios}
\end{table}

El criterio SC4, la calidad visual de producto, se cerró con la aprobación del autor. Empezó
como aprobación anticipada condicional, mientras el autor estaba desconectado. Quedó
confirmada sin condiciones tras una revisión remota de dieciocho capturas del entorno local
que cubrían las cuatro superficies principales en escritorio y en móvil, ambos temas y una
segunda cuenta simulada con colección vacía como prueba visual de independencia.

\section{Revisión de código de todo el repositorio}\label{sec:pru-revision}

Al cierre de la Fase 01.1, un subagente revisor examinó todo el árbol con una postura
adversarial de búsqueda de defectos y un contexto propio, distinto del de la implementación.
La revisión encontró dieciséis hallazgos, cuya distribución por severidad recoge la
Tabla~\ref{tab:hallazgos}. El autor autorizó aplicar la totalidad de los hallazgos, que se
corrigieron en cinco tandas confirmadas, cada una con las pruebas ejecutadas después.

\begin{table}[htp]
\centering
\begin{tabular}{ | l | c | c | }
    \hline
    \textbf{Severidad} & \textbf{Hallazgos} & \textbf{Estado} \\
    \hline
    Crítica & 0 & Ninguno \\
    \hline
    Alta & 3 & Corregidos \\
    \hline
    Media & 6 & Corregidos \\
    \hline
    Baja & 7 & Corregidos \\
    \hline
    \textbf{Total} & \textbf{16} & \textbf{16 corregidos} \\
    \hline
\end{tabular}
\caption{Hallazgos de la revisión de código de todo el repositorio.}
\label{tab:hallazgos}
\end{table}

Los tres hallazgos de severidad alta ilustran el valor de una revisión con contexto
independiente. El primero era que la configuración de DRF dejaba habilitada de forma
silenciosa la autenticación básica en todo punto de acceso autenticado, lo que abría una vía
de adivinación de contraseñas sin límite de ritmo y una ruta de cambio de estado exenta de la
protección contra falsificación de petición en sitios cruzados. El segundo era que una
reimportación troceada del catálogo de IGDB tras una pasada completa podía saltarse en
silencio un rango de identificadores ya recorrido y aun así declararse completa con una huella
nueva. El tercero era que el arranque de producción usaba el servidor de desarrollo de
Django, que su documentación desaconseja explícitamente. Los tres se corrigieron, el segundo
con un campo de esquema nuevo y una migración, el tercero con la incorporación del servidor
\texttt{gunicorn}.

La revisión también dejó constancia de las áreas sanas: ausencia de inyección de SQL, con
todo el uso del ORM parametrizado; validación minuciosa de las consultas del catálogo;
acotación de propiedad en todas las vistas de datos personales; escrituras transaccionales
con bloqueo de fila y restricciones de unicidad como guarda real de carreras; higiene de
secretos con redacción en registros y excepciones.

\section{Lectura de los resultados}\label{sec:pru-lectura}

Las cifras anteriores respaldan afirmaciones acotadas. El sistema tiene una cobertura de
pruebas amplia sobre PostgreSQL real, una accesibilidad automática sólida aunque no
exhaustiva y una revisión de seguridad de todo el árbol con sus hallazgos corregidos. Esas
cifras no respaldan ninguna medida de rendimiento bajo carga, de recuperación ante desastre
ni de calidad de recomendación. La memoria tampoco las afirma, porque pertenecen a fases
posteriores. Los dos elementos manuales de accesibilidad de la Fase 1 siguen abiertos y
declarados como tales.

Este capítulo ha mostrado cómo se verifica el sistema. El capítulo siguiente trata su
despliegue.
